\documentclass[10pt,a4paper]{article}
\usepackage[utf8]{inputenc}
\usepackage{amsmath, amsthm, amssymb}
\usepackage{geometry}
\usepackage{enumitem}
\usepackage{bm}

\geometry{
  a4paper,
  left=20mm,
  right=20mm,
  top=25mm,
  bottom=25mm,
}

\begin{document}
\section*{Problem}
\begin{enumerate}[label=\fbox{\arabic*}]
  \item % Problem 1
        Let $p$ be a prime number greater than or equal to $3$. We denote by $\mathrm{SL}(2, \mathbb{F}_p)$ the group of all $2 \times 2$-matrices $A$ whose entries consist of elements of the finite field $\mathbb{F}_p$ and such that $\det(A) = 1$. Determine the number of elements $A$ of $\mathrm{SL}(2, \mathbb{F}_p)$ such that $A^{p-1} = \mathbb{I}$, where $\mathbb{I}$ is the unit matrix.

  \item % Problem 2
        Let $K = \mathbb{C}(t)$ be the field of rational functions over $\mathbb{C}$ in variable $t$. Let $u$ be a nonzero complex number and let $L$ be the splitting field of the polynomial $f(X) = X^4 + 2utX^2 + t \in K[X]$ over $K$.
        \begin{enumerate}[label=(\arabic*)]
          \item Determine the degree $[L : K]$.
          \item Is the Galois group $\mathrm{Gal}(L/K)$ abelian? Give a reason for your answer.
        \end{enumerate}

  \item % Problem 3
        Let $m$ and $n$ be integers greater than or equal to $2$, and let $a_1, \dots, a_n$ be integers with $0 \le a_i < m$. Let $A = \mathbb{C}[x_1, \dots, x_n]$ be a polynomial ring over $\mathbb{C}$ with variables $x_1, \dots, x_n$. To each $P \in A$, we associate $\varphi(P) \in A$ by
        \[
          \varphi(P)(x_1, \dots, x_n) = P(\zeta_m^{a_1} x_1, \dots, \zeta_m^{a_n} x_n),
        \]
        where $\zeta_m := e^{2\pi\sqrt{-1}/m}$ is the $m$-th primitive root and define $A' := \{ P \in A \mid P = \varphi(P) \}$.
        \begin{enumerate}[label=(\arabic*)]
          \item Prove that $A'$ is a subring of $A$.
          \item We suppose that $a_i$ are all at most $1$. Under this condition, find a necessary and sufficient condition such that $A'$ is isomorphic to a polynomial ring over $\mathbb{C}$.
        \end{enumerate}

  \item % Problem 4
        Let $\mathbb{RP}^n$ denote the $n$-dimensional real projective space. For $(x_0, x_1, \dots, x_n) \in \mathbb{R}^{n+1} \setminus \{\mathbf{0}\}$, let $[x_0, x_1, \dots, x_n] \in \mathbb{RP}^n$ denote the point that is given by the $1$-dimensional real vector subspace generated by $(x_0, x_1, \dots, x_n)$. Define a subset $M$ of $\mathbb{RP}^3$ by
        \[
          M = \{ [x_0, x_1, x_2, x_3] \in \mathbb{RP}^3 \mid x_0^2 + x_1^2 = x_2^2 + x_3^2 \}.
        \]
        \begin{enumerate}[label=(\arabic*)]
          \item Show that $M$ is a $C^\infty$-submanifold of $\mathbb{RP}^3$.
          \item Let $f : M \to \mathbb{RP}^1 \times \mathbb{RP}^1$ be the map defined by
          \[
            f([x_0, x_1, x_2, x_3]) = ([x_0, x_1], [x_2, x_3]).
          \]
          Show that $f$ is a $C^\infty$-submersion.
          \item Let $S^1 = \{ (x, y) \in \mathbb{R}^2 \mid x^2 + y^2 = 1 \}$. Show that $M$ is diffeomorphic to $S^1 \times S^1$.
        \end{enumerate}

  \item % Problem 5
        Let $D^2 = \{ (x, y) \in \mathbb{R}^2 \mid x^2 + y^2 \le 1 \}$ and $S^1 = \{ (x, y) \in \mathbb{R}^2 \mid x^2 + y^2 = 1 \}$. Define subspaces $X_1, X_2, X_3$ of $D^2 \times D^2 \times D^2$ by
        \[
          X_1 = D^2 \times S^1 \times S^1, \quad X_2 = S^1 \times D^2 \times S^1, \quad X_3 = S^1 \times S^1 \times D^2.
        \]
        \begin{enumerate}[label=(\arabic*)]
          \item Compute the homology groups of $Y = X_1 \cup X_2$ with integer coefficients.
          \item Compute the homology groups of $Z = X_1 \cup X_2 \cup X_3$ with integer coefficients.
        \end{enumerate}

  \item % Problem 6
        Let $\mu$ be a Borel measure on $[0, \infty)$. Let $f$ be a function on $[0, \infty)$ which is uniformly continuous and non-negative. Assume
        \[
          \int_{[0, \infty)} e^{f(x)} \mu(dx) < \infty.
        \]
        Find the following limit:
        \[
          \lim_{n \to \infty} \int_{[0, \infty)} \left( 1 + \int_0^{1/n} f(x+t) dt \right)^n \mu(dx).
        \]

  \item % Problem 7
        Let $I = (0, 1)$ and let $L^2(I)$ be the space of all square integrable complex-valued functions on $I$. We set
        \[
          (f, g) := \int_I f(x) \overline{g(x)} dx, \quad \|f\|_2 := \sqrt{(f, f)}, \quad \|f\|_1 := \int_I |f(x)| dx
        \]
        Answer the following questions.
        \begin{enumerate}[label=(\arabic*)]
          \item Suppose that a sequence $\{f_n\}_{n=1}^\infty$ in $L^2(I)$ satisfies
          \[
            \|f_n\|_1 \longrightarrow 0, \quad \|f_n\|_2 = 1 \ (n \ge 1).
          \]
          Show that $f_n \longrightarrow 0$ weakly in $L^2(I)$.
          \item Suppose that $T : L^2(I) \to L^2(I)$ is compact. Show that
          \[
            \forall \varepsilon > 0, \ \exists C_\varepsilon > 0; \ \forall f \in L^2(I), \ \|Tf\|_2 \le \varepsilon \|f\|_2 + C_\varepsilon \|f\|_1.
          \]
        \end{enumerate}

  \item % Problem 8
        Let $n$ be an integer with $n \ge 2$, $\Omega = \{ x \in \mathbb{R}^n : |x| < 1 \}$, and $\overline{\Omega}$ be its closure. Let $C^2(\overline{\Omega})$ be the set of all real valued $C^2$ functions on $\Omega$ such that all partial derivatives up to order $2$ can be extended to $\overline{\Omega}$ continuously. Assume that $f \in C^2(\overline{\Omega})$. For any $\varepsilon > 0$, let $u_\varepsilon \in C^2(\overline{\Omega})$ be a solution of
        \begin{align*}
          -\varepsilon \Delta u_\varepsilon + u_\varepsilon &= f, \quad x \in \Omega, \\
          u_\varepsilon &= 0, \quad x \in \partial\Omega,
        \end{align*}
        where $\Delta = \sum_{i=1}^n \frac{\partial^2}{\partial x_i^2}$.
        \begin{enumerate}[label=(\roman*)]
          \item Show that there exist $C_1, C_2 > 0$ independent of $\varepsilon$ such that
          \[
            \sqrt{\varepsilon} \|\nabla u_\varepsilon\|_{L^2(\Omega)} \le C_1 \|f\|_{L^2(\Omega)} \quad \text{and} \quad \varepsilon \|\Delta u_\varepsilon\|_{L^2(\Omega)} \le C_2 \|f\|_{L^2(\Omega)}.
          \]
          \item Prove that $\lim_{\varepsilon \downarrow 0} \|u_\varepsilon - f\|_{L^2(\Omega)} = 0$.
        \end{enumerate}

  \item % Problem 9
        Define the set $X$ of real-valued functions by
        \[
          X = \left\{ u \in C^2((0, 1)) \cap C([0, 1]) \ \Big| \ u(0) = 0, \ u(1) = 1, \ \int_0^1 (x u'(x))^2 dx < \infty \right\}
        \]
        and for $u \in X$ consider the following functional:
        \[
          J[u] = \int_0^1 e^{-x} \left( 2(u(x))^2 + \frac{9}{2}x^2 (u'(x))^2 \right) dx.
        \]
        Answer the following questions.
        \begin{enumerate}[label=(\arabic*)]
          \item Let $C_0^\infty((0, 1))$ denote the space of all real-valued infinitely differentiable functions with support in the open interval $(0, 1)$. For $u \in X$ and $\varphi \in C_0^\infty((0, 1))$ compute
          \[
            \frac{\delta J}{\delta u}(u; \varphi) := \lim_{\varepsilon \to 0} \frac{J[u + \varepsilon \varphi] - J[u]}{\varepsilon}.
          \]
          \item Find in series form all $u \in X$ that satisfy the following condition:
          \[
            \frac{\delta J}{\delta u}(u; \varphi) = 0 \quad \text{whenever} \quad \varphi \in C_0^\infty((0, 1)).
          \]
        \end{enumerate}

  \item % Problem 10
        Let $\mathbf{N}$ be the set of nonnegative integers and $B = \{ x_1 \dots x_n \mid x_1, \dots, x_n \in \{\mathbf{a}, \mathbf{b}\}, n \ge 0 \}$ be the set of finite strings of characters $\mathbf{a}$ and $\mathbf{b}$. Let us write $\epsilon$ to denote the null string and $x^m$ to denote the string of $m$ copies of the character $x$. Given a string $w \in B$, we write $|w|$ for the length of $w$ and $|w|_\mathbf{a}$ for the number of occurrences of $\mathbf{a}$ in $w$.

        Let us inductively define a function $g : B \times B \to B$ as follows.
        \begin{align*}
          g(\epsilon, u) &= u \\
          g(\mathbf{a}w, u) &= g(w, u\mathbf{a}) \\
          g(\mathbf{b}\mathbf{a}w, u) &= u\mathbf{a}\mathbf{b}w \\
          g(\mathbf{b}\mathbf{b}w, u) &= g(\mathbf{b}w, \mathbf{b}u) \\
          g(\mathbf{b}, u) &= \mathbf{b}u
        \end{align*}
        Also, define a function $f : B \to B$ by $f(w) = g(w, \epsilon)$.

        Suppose $1 \le k < n$. Letting $B_{n,k} = \{ w \in B \mid |w| = n, |w|_\mathbf{a} = k \}$, show that the following equation holds.
        \[
          \{ f^i(\mathbf{b}^{n-k}\mathbf{a}^k) \mid i \in \mathbf{N} \} = B_{n,k}
        \]
\end{enumerate}

\newpage
\section*{Solution}
\begin{enumerate}[label=\fbox{\arabic*}]
  \item % Solution 1
        For $A \in \mathrm{SL}(2, \mathbb{F}_p)$, its characteristic polynomial is $x^2 - \mathrm{tr}(A)x + 1 = 0$. The roots of this polynomial are the eigenvalues of $A$, let them be $\lambda$ and $\lambda^{-1}$. These eigenvalues lie either in $\mathbb{F}_p$ or in the quadratic extension $\mathbb{F}_{p^2}$. We classify the elements $A$ based on their eigenvalues:
        \begin{itemize}
            \item \textbf{Case 1: Eigenvalues are in $\mathbb{F}_{p^2} \setminus \mathbb{F}_p$.}
            The eigenvalues are of the form $\lambda$ and $\lambda^p$, with $\lambda \cdot \lambda^p = \lambda^{p+1} = 1$. If $A^{p-1} = \mathbb{I}$, its eigenvalues must satisfy $\lambda^{p-1} = 1$. However, $\lambda^{p+1} = 1$ and $\lambda^{p-1} = 1$ implies $\lambda^2 = 1$, so $\lambda = \pm 1 \in \mathbb{F}_p$, which contradicts that the eigenvalues are strictly in $\mathbb{F}_{p^2} \setminus \mathbb{F}_p$. Thus, there are $0$ such matrices satisfying the condition.
            
            \item \textbf{Case 2: Eigenvalues are in $\mathbb{F}_p$.}
            Since $\lambda \in \mathbb{F}_p^\times$, by Fermat's Little Theorem, $\lambda^{p-1} = 1$. Thus, \emph{any diagonalizable} matrix over $\mathbb{F}_p$ satisfies $A^{p-1} = \mathbb{I}$.
            If $\lambda \neq \pm 1$, there are $\frac{p-3}{2}$ distinct pairs of eigenvalues $\{\lambda, \lambda^{-1}\}$. For each pair, the centralizer is the set of diagonal matrices of size $p-1$. The number of matrices in a conjugacy class for a fixed pair is exactly $\frac{|\mathrm{SL}(2, \mathbb{F}_p)|}{p-1} = \frac{p(p-1)(p+1)}{p-1} = p(p+1)$. This gives $\frac{p-3}{2} p(p+1)$ matrices.
            
            If $\lambda = 1$ or $\lambda = -1$, the centralizer of the non-diagonalizable matrices (with Jordan blocks) differs, but we can test them directly. A non-diagonalizable matrix with repeated eigenvalue $\pm 1$ is similar to $J = \pm \begin{pmatrix} 1 & 1 \\ 0 & 1 \end{pmatrix}$. Since $p \ge 3$, $p-1$ is even, so $J^{p-1} = \begin{pmatrix} 1 & p-1 \\ 0 & 1 \end{pmatrix} \neq \mathbb{I}$ (as $p-1 \not\equiv 0 \pmod p$). Thus, no non-diagonalizable matrices satisfy $A^{p-1} = \mathbb{I}$. 
            The only remaining matrices are the purely scalar ones, $A = \mathbb{I}$ and $A = -\mathbb{I}$, both of which trivially satisfy the equation (since $p-1$ is even). This gives $2$ matrices.
        \end{itemize}
        Summing these up, the total number of elements satisfying $A^{p-1} = \mathbb{I}$ is:
        \[
          2 + \frac{p(p+1)(p-3)}{2} = \frac{p^3 - 2p^2 - 3p + 4}{2}.
        \]

  \item % Solution 2
        \begin{enumerate}[label=(\arabic*)]
          \item Let $Y = X^2$. The polynomial becomes $g(Y) = Y^2 + 2utY + t = 0$. Using the quadratic formula, the roots are $Y = -ut \pm \sqrt{u^2t^2 - t}$. By Eisenstein's criterion on $\mathbb{C}[t]$ at the prime ideal $(t)$, $f(X)$ is irreducible over $K$. 
          The four roots of $f(X)$ are $\pm \alpha, \pm \beta$ where $\alpha^2 = -ut + \sqrt{t(u^2t-1)}$ and $\beta^2 = -ut - \sqrt{t(u^2t-1)}$. The splitting field is $L = K(\alpha, \beta)$. Let $s = \sqrt{t}$. Then $\alpha^2 \beta^2 = t \implies \beta = s/\alpha$. Thus, $L$ contains $s = \alpha\beta$. 
          Using the denesting formula for square roots, $\alpha = \sqrt{A} + \sqrt{B}$ where $A = \frac{s - us^2}{2}$ and $B = \frac{-s - us^2}{2}$. We can express the roots as linear combinations of $\sqrt{A}$ and $\sqrt{B}$. Since $AB = \frac{s^2(u^2s^2 - 1)}{4}$ is not a square in $K(s)$, the field extension adjoining both $\sqrt{A}$ and $\sqrt{B}$ over $K(s)$ has degree $4$. Since $[K(s) : K] = 2$, the total degree is $[L : K] = 2 \times 4 = 8$.

          \item \textbf{No, the Galois group is not abelian.} 
          The Galois group of an irreducible biquadratic polynomial $X^4 + aX^2 + b$ over $K$ is completely determined by whether $b$, $a^2-4b$, or $b(a^2-4b)$ are perfect squares in $K$. Here, $b = t$, $a^2-4b = 4t(u^2t-1)$, and $b(a^2-4b) = 4t^2(u^2t-1)$. None of these expressions are perfect squares in $\mathbb{C}(t)$. Thus, the Galois group $\mathrm{Gal}(L/K)$ is isomorphic to the dihedral group of order $8$ ($D_4$), which is well-known to be non-abelian.
        \end{enumerate}

  \item % Solution 3
        \begin{enumerate}[label=(\arabic*)]
          \item We verify the ring axioms for $A'$. Since $\varphi$ acts purely by scaling the variables by roots of unity, it is a $\mathbb{C}$-algebra homomorphism. For any $P, Q \in A'$, we have $\varphi(P - Q) = \varphi(P) - \varphi(Q) = P - Q$, meaning $P - Q \in A'$. Furthermore, $\varphi(PQ) = \varphi(P)\varphi(Q) = PQ$, meaning $PQ \in A'$. Finally, $\varphi(1) = 1$, so the multiplicative identity is in $A'$. Therefore, $A'$ is a subring of $A$.

          \item Because $0 \le a_i \le 1$, the transformation $\varphi$ acts as the identity on variables $x_i$ where $a_i = 0$, and multiplies $x_i$ by $\zeta_m$ where $a_i = 1$. Let $k$ be the number of indices where $a_i = 1$. By the Chevalley-Shephard-Todd theorem, the invariant ring of a finite group acting linearly on a polynomial ring is isomorphic to a polynomial ring if and only if the group is generated by pseudo-reflections. 
          The action of $\varphi$ defines a cyclic group $C_m$. The generator $\varphi$ acts as a pseudo-reflection (an element fixing a hyperplane pointwise) if and only if exactly one eigenvalue is non-trivially a root of unity, meaning $k = 1$ (or $k=0$ for the identity action). 
          If $k \ge 2$, there are at least two eigenvalues equal to $\zeta_m$, meaning the space fixed by $\varphi$ has codimension $\ge 2$. Since $m \ge 2$, the group cannot be generated by pseudo-reflections, making $A'$ not isomorphic to a polynomial ring. 
          Therefore, the necessary and sufficient condition is that $k \le 1$, which translates to:
          \[
            \sum_{i=1}^n a_i \le 1.
          \]
        \end{enumerate}

  \item % Solution 4
        \begin{enumerate}[label=(\arabic*)]
          \item Define the homogeneous polynomial $F(x) = x_0^2 + x_1^2 - x_2^2 - x_3^2$. The set $M$ is the projective zero locus of $F$. In $\mathbb{R}^4$, the gradient $\nabla F(x) = (2x_0, 2x_1, -2x_2, -2x_3)$ vanishes only at the origin $\mathbf{0}$. Since $\mathbf{0}$ is not in the projective space, $0$ is a regular value of $F$ on the sphere $S^3$. By the Regular Level Set Theorem, the projection to $\mathbb{RP}^3$ yields a smooth $C^\infty$-submanifold of dimension $2$.

          \item The map $f$ is given by $f([x]) = ([x_0, x_1], [x_2, x_3])$. Since points in $M$ satisfy $x_0^2+x_1^2 = x_2^2+x_3^2$, and $x \neq \mathbf{0}$, neither $(x_0, x_1)$ nor $(x_2, x_3)$ can be the zero vector. Thus $f$ is well-defined and manifestly smooth.
          Both $M$ and $\mathbb{RP}^1 \times \mathbb{RP}^1$ have dimension $2$. We construct the inverse locally: for any $([u], [v]) \in \mathbb{RP}^1 \times \mathbb{RP}^1$, normalize them to unit length $u_0^2+u_1^2=1$ and $v_0^2+v_1^2=1$. The point $[u_0, u_1, v_0, v_1]$ lies in $M$. Negating $u$ or $v$ independently gives $4$ choices in $\mathbb{R}^4$, condensing to $2$ antipodal pairs in $\mathbb{RP}^3$. Thus, $f$ is a $2$-to-$1$ smooth covering map. Since covering maps are local diffeomorphisms, its differential is everywhere invertible, making $f$ a $C^\infty$-submersion.

          \item A point $[x] \in M$ can be scaled by some $\lambda \neq 0$. To analyze the topology, we can restrict our homogeneous coordinates to the subspace where $x_0^2 + x_1^2 = x_2^2 + x_3^2 = 1$. This subspace in $\mathbb{R}^4$ is exactly $S^1 \times S^1$.
          The equivalence relation from the projective scaling requires $\lambda^2 = 1 \implies \lambda = \pm 1$. Therefore, $M = (S^1 \times S^1) / \sim$ where $(u, v) \sim (-u, -v)$. 
          We can map this quotient diffeomorphically to $S^1 \times S^1$ (viewed in $\mathbb{C}^2$) via the map $g(u, v) = (u^2, uv)$. Note that $g(-u, -v) = ((-u)^2, (-u)(-v)) = (u^2, uv)$, so it is well-defined on the quotient. It possesses a smooth inverse locally (by taking the branches of the square root $u = \pm\sqrt{z}$ and $v = w/u$), making it a global diffeomorphism. Thus $M \cong S^1 \times S^1$.
        \end{enumerate}

  \item % Solution 5
        \begin{enumerate}[label=(\arabic*)]
          \item Let $Y = X_1 \cup X_2$. Note that $X_1 \cong D^2 \times T^2$ and $X_2 \cong D^2 \times T^2$. Since they share the $S^1$ factor in the third coordinate, we can rewrite their union structurally as $(D^2 \times S^1 \cup S^1 \times D^2) \times S^1$. The first term is a known decomposition of the 3-sphere $S^3$ formed by joining two solid tori along their boundaries. Thus, $Y \cong S^3 \times S^1$.
          The homology groups of $S^3 \times S^1$ can be directly computed via the Künneth formula (or verified through Mayer-Vietoris on $X_1$ and $X_2$ intersecting at $T^3$). The homology groups with integer coefficients are:
          \[
            H_0(Y) = \mathbb{Z}, \quad H_1(Y) = \mathbb{Z}, \quad H_2(Y) = 0, \quad H_3(Y) = \mathbb{Z}, \quad H_4(Y) = \mathbb{Z}.
          \]

          \item Now we compute the homology of $Z = Y \cup X_3 = (S^3 \times S^1) \cup (S^1 \times S^1 \times D^2)$. The intersection is $Y \cap X_3 = (X_1 \cap X_3) \cup (X_2 \cap X_3) = T^3 \cup T^3 = T^3$, matching the boundary of $X_3$.
          We apply the Mayer-Vietoris sequence for $Z = Y \cup X_3$:
          \[
            \dots \to H_k(T^3) \to H_k(S^3 \times S^1) \oplus H_k(T^2 \times D^2) \to H_k(Z) \to H_{k-1}(T^3) \to \dots
          \]
          For $k=4$: $H_4(Y) = \mathbb{Z}$ and $H_3(T^3) = \mathbb{Z}$. The generator of $H_3(T^3)$ projects to the Clifford torus times $S^1$ in $S^3 \times S^1$, which represents a trivial homology class in $S^3$ (since $H_3(T^3) \to H_3(Y)$ is zero). Thus $H_4(Z) = \mathbb{Z} \oplus \mathbb{Z} = \mathbb{Z}^2$.
          For $k=3$: $H_3(T^3) \to 0$ means we have $0 \to \mathbb{Z} \to H_3(Z) \to \mathbb{Z}^3 \xrightarrow{\psi_2} 0 \oplus \mathbb{Z}$. The map $\psi_2$ has a 2-dimensional kernel, yielding $H_3(Z) \cong \mathbb{Z} \oplus \mathbb{Z}^2 = \mathbb{Z}^3$.
          For $k=2$: $H_1(T^3) \xrightarrow{\psi_1} H_1(Y) \oplus H_1(X_3)$ maps $\mathbb{Z}^3 \to \mathbb{Z} \oplus \mathbb{Z}^2$ as an isomorphism. Hence $H_2(Z) = 0$.
          For $k=1$: Because $\psi_1$ is an isomorphism, its cokernel is trivial, giving $H_1(Z) = 0$.
          Summarizing the homology groups of $Z$:
          \[
            H_0(Z) = \mathbb{Z}, \quad H_1(Z) = 0, \quad H_2(Z) = 0, \quad H_3(Z) = \mathbb{Z}^3, \quad H_4(Z) = \mathbb{Z}^2.
          \]
        \end{enumerate}

  \item % Solution 6
        Let $g_n(x) = \left( 1 + \int_0^{1/n} f(x+t) dt \right)^n$. We define the average $\bar{f}_n(x) = n \int_0^{1/n} f(x+t) dt$, allowing us to write $g_n(x) = \left( 1 + \frac{\bar{f}_n(x)}{n} \right)^n$.
        Since $f$ is continuous, by the Fundamental Theorem of Calculus, $\lim_{n \to \infty} \bar{f}_n(x) = f(x)$. Thus, the integrand $g_n(x)$ converges pointwise to $e^{f(x)}$.
        To justify exchanging the limit and the integral, we use the Lebesgue Dominated Convergence Theorem (DCT). Notice that for any $y \ge 0$, $(1 + y/n)^n \le e^y$, so $g_n(x) \le e^{\bar{f}_n(x)}$. 
        Because $f$ is uniformly continuous on $[0, \infty)$, there exists $\delta > 0$ such that $|t| < \delta \implies |f(x+t) - f(x)| < 1$. For large enough $n$ (specifically $n > 1/\delta$), $t \in [0, 1/n]$ implies $f(x+t) < f(x) + 1$. Integrating this over the interval yields $\bar{f}_n(x) \le f(x) + 1$.
        This gives a uniform bounding function:
        \[
          g_n(x) \le e^{f(x)+1} = e \cdot e^{f(x)}.
        \]
        We are given that $\int_{[0, \infty)} e^{f(x)} \mu(dx) < \infty$, so the dominating function $e \cdot e^{f(x)}$ is integrable with respect to $\mu$. By the DCT, we can pass the limit inside the integral:
        \[
          \lim_{n \to \infty} \int_{[0, \infty)} g_n(x) \mu(dx) = \int_{[0, \infty)} \lim_{n \to \infty} g_n(x) \mu(dx) = \int_{[0, \infty)} e^{f(x)} \mu(dx).
        \]

  \item % Solution 7
        \begin{enumerate}[label=(\arabic*)]
          \item To show $f_n \to 0$ weakly in $L^2(I)$, we must show that $\langle f_n, g \rangle \to 0$ for all $g \in L^2(I)$.
          Because $\|f_n\|_2 = 1$, the sequence is bounded in the Hilbert space $L^2(I)$. By a standard density argument, it suffices to verify weak convergence on a dense subset. We choose $L^\infty(I)$, which is dense in $L^2(I)$ (since $I = (0, 1)$ has finite measure).
          Let $g \in L^\infty(I)$. We can bound the inner product as follows:
          \[
            |\langle f_n, g \rangle| \le \int_0^1 |f_n(x) \overline{g(x)}| dx \le \|g\|_\infty \int_0^1 |f_n(x)| dx = \|g\|_\infty \|f_n\|_1.
          \]
          Since we are given that $\|f_n\|_1 \to 0$, it follows that $|\langle f_n, g \rangle| \to 0$. By density, this extends to all $g \in L^2(I)$, confirming that $f_n$ converges weakly to $0$.

          \item We proceed by contradiction. Suppose the statement is false. Then there exists an $\varepsilon > 0$ such that for all $n \ge 1$, we can find a function $h_n \in L^2(I)$ where
          \[
            \|Th_n\|_2 > \varepsilon \|h_n\|_2 + n \|h_n\|_1.
          \]
          By scaling $f_n = h_n / \|h_n\|_2$, we can normalize the sequence such that $\|f_n\|_2 = 1$. The inequality then becomes:
          \[
            \|Tf_n\|_2 > \varepsilon + n \|f_n\|_1.
          \]
          Since $T$ is a bounded operator (as all compact operators are), $\|Tf_n\|_2 \le \|T\| \|f_n\|_2 = \|T\|$. Consequently,
          \[
            n \|f_n\|_1 < \|Tf_n\|_2 - \varepsilon \le \|T\| \implies \|f_n\|_1 < \frac{\|T\|}{n}.
          \]
          Taking the limit as $n \to \infty$ shows that $\|f_n\|_1 \to 0$. From part (1), since $\|f_n\|_1 \to 0$ and $\|f_n\|_2 = 1$, we know $f_n \to 0$ weakly in $L^2(I)$.
          A fundamental property of compact operators on Hilbert spaces is that they map weakly convergent sequences to strongly convergent sequences. Therefore, since $f_n \rightharpoonup 0$, we must have $Tf_n \to T(0) = 0$ strongly in $L^2(I)$, which implies $\|Tf_n\|_2 \to 0$.
          However, our initial hypothesis gave $\|Tf_n\|_2 > \varepsilon > 0$ for all $n$, representing a direct contradiction. Thus, the original claim must be true.
        \end{enumerate}

  \item % Solution 8
        \begin{enumerate}[label=(\roman*)]
          \item We multiply the PDE $-\varepsilon \Delta u_\varepsilon + u_\varepsilon = f$ by $u_\varepsilon$ and integrate over $\Omega$. Using Green's first identity (integration by parts) and the Dirichlet boundary condition $u_\varepsilon = 0$ on $\partial\Omega$:
          \[
            -\varepsilon \int_\Omega u_\varepsilon \Delta u_\varepsilon dx + \int_\Omega u_\varepsilon^2 dx = \varepsilon \int_\Omega |\nabla u_\varepsilon|^2 dx + \int_\Omega u_\varepsilon^2 dx = \int_\Omega f u_\varepsilon dx.
          \]
          Applying the Cauchy-Schwarz inequality to the right hand side gives $\int f u_\varepsilon \le \|f\|_{L^2} \|u_\varepsilon\|_{L^2}$. This immediately implies two bounds:
          1) $\|u_\varepsilon\|_{L^2}^2 \le \|f\|_{L^2} \|u_\varepsilon\|_{L^2} \implies \|u_\varepsilon\|_{L^2} \le \|f\|_{L^2}$.
          2) $\varepsilon \|\nabla u_\varepsilon\|_{L^2}^2 \le \|f\|_{L^2} \|u_\varepsilon\|_{L^2} \le \|f\|_{L^2}^2 \implies \sqrt{\varepsilon} \|\nabla u_\varepsilon\|_{L^2} \le \|f\|_{L^2}$. (Thus $C_1 = 1$).
          
          To find the bound for the Laplacian, we isolate it algebraically: $-\varepsilon \Delta u_\varepsilon = f - u_\varepsilon$.
          Taking the $L^2$ norm of both sides and using the triangle inequality with the $L^2$ bound on $u_\varepsilon$:
          \[
            \varepsilon \|\Delta u_\varepsilon\|_{L^2} = \|f - u_\varepsilon\|_{L^2} \le \|f\|_{L^2} + \|u_\varepsilon\|_{L^2} \le 2\|f\|_{L^2}.
          \]
          (Alternatively, multiplying by $-\Delta u_\varepsilon$ and integrating yields $C_2 = 1$). Thus the bounds hold independent of $\varepsilon$.

          \item We rewrite the PDE as $u_\varepsilon - f = \varepsilon \Delta u_\varepsilon$. Our goal is to show $\|u_\varepsilon - f\|_{L^2} \to 0$.
          We utilize a density argument. Let $\delta > 0$ be given. Since $C_c^\infty(\Omega)$ is dense in $L^2(\Omega)$, there exists a function $g \in C_c^\infty(\Omega)$ such that $\|f - g\|_{L^2} < \delta$. Let $h = f - g$, so $f = g + h$.
          By linearity of the PDE, the solution $u_\varepsilon(f)$ corresponding to $f$ can be decomposed as $u_\varepsilon(g) + u_\varepsilon(h)$. By the triangle inequality:
          \[
            \|u_\varepsilon(f) - f\|_{L^2} \le \|u_\varepsilon(g) - g\|_{L^2} + \|u_\varepsilon(h) - h\|_{L^2}.
          \]
          For the $h$ term, from part (i) we know $\|u_\varepsilon(h)\|_{L^2} \le \|h\|_{L^2}$. Thus:
          \[
            \|u_\varepsilon(h) - h\|_{L^2} \le \|u_\varepsilon(h)\|_{L^2} + \|h\|_{L^2} \le 2\|h\|_{L^2} < 2\delta.
          \]
          For the smooth, compactly supported $g$, let $w = u_\varepsilon(g) - g$. Substituting into the PDE gives $-\varepsilon \Delta w + w = \varepsilon \Delta g$. Since $g$ vanishes on $\partial\Omega$, $w \in H_0^1(\Omega)$. Multiplying by $w$ and integrating yields:
          \[
            \varepsilon \|\nabla w\|_{L^2}^2 + \|w\|_{L^2}^2 = \varepsilon \int_\Omega (\Delta g) w dx \le \varepsilon \|\Delta g\|_{L^2} \|w\|_{L^2}.
          \]
          This implies $\|w\|_{L^2} \le \varepsilon \|\Delta g\|_{L^2}$. As $\varepsilon \downarrow 0$, $\|u_\varepsilon(g) - g\|_{L^2} \to 0$.
          Taking the limit superior of the whole expression gives $\limsup_{\varepsilon \downarrow 0} \|u_\varepsilon(f) - f\|_{L^2} \le 2\delta$. Since $\delta > 0$ is arbitrary, the limit is strictly $0$.
        \end{enumerate}

  \item % Solution 9
        \begin{enumerate}[label=(\arabic*)]
          \item We substitute $u + \varepsilon\varphi$ into the functional $J$:
          \[
            J[u + \varepsilon\varphi] = \int_0^1 e^{-x} \left( 2(u + \varepsilon\varphi)^2 + \frac{9}{2}x^2 (u' + \varepsilon\varphi')^2 \right) dx.
          \]
          Expanding this and differentiating with respect to $\varepsilon$ at $\varepsilon = 0$ provides the required limit (the Fréchet derivative):
          \[
            \frac{\delta J}{\delta u}(u; \varphi) = \int_0^1 e^{-x} \left( 4u\varphi + 9x^2 u' \varphi' \right) dx.
          \]
          We apply integration by parts to the term containing $\varphi'$ to shift the derivative onto the other factors:
          \[
            \int_0^1 (9x^2 e^{-x} u') \varphi' dx = \Big[ 9x^2 e^{-x} u' \varphi \Big]_0^1 - \int_0^1 (9x^2 e^{-x} u')' \varphi dx.
          \]
          Because $\varphi \in C_0^\infty((0, 1))$, it vanishes at the boundaries $x=0$ and $x=1$, rendering the boundary evaluation zero. Thus, the variation evaluates to:
          \[
            \frac{\delta J}{\delta u}(u; \varphi) = \int_0^1 \left( 4e^{-x}u - (9x^2 e^{-x} u')' \right) \varphi dx.
          \]

          \item Setting the variation to zero for all $\varphi \in C_0^\infty((0, 1))$ implies the integrand must vanish (Euler-Lagrange equation):
          \[
            4e^{-x}u - (9x^2 e^{-x} u')' = 0 \implies 4e^{-x}u - 18x e^{-x}u' + 9x^2 e^{-x}u' - 9x^2 e^{-x} u'' = 0.
          \]
          Dividing out the common $e^{-x}$ term simplifies the ODE to:
          \[
            9x^2 u'' + 9x(2 - x) u' - 4u = 0.
          \]
          We seek a series solution via the Frobenius method: $u(x) = \sum_{k=0}^\infty a_k x^{k+r}$ with $a_0 \neq 0$.
          Substituting $u$, $u'$, and $u''$ into the ODE gives:
          \[
            \sum_{k=0}^\infty \big[ 9(k+r)(k+r-1) + 18(k+r) - 4 \big] a_k x^{k+r} - \sum_{k=0}^\infty 9(k+r) a_k x^{k+r+1} = 0.
          \]
          For $k=0$, the indicial equation is $9r(r-1) + 18r - 4 = 0$, which factors as $(3r-1)(3r+4) = 0$. Since $u(0) = 0$, we require the positive root $r = 1/3$.
          Shifting the index on the second sum and equating coefficients for $k \ge 1$ with $r=1/3$ gives:
          \[
            a_k \big[ 9(k+1/3)(k-2/3) + 18(k+1/3) - 4 \big] = 9(k-2/3) a_{k-1}.
          \]
          Simplifying the bracket on the LHS reduces elegantly to $3k(3k+5)$. The recurrence relation is therefore:
          \[
            3k(3k+5) a_k = 3(3k-2) a_{k-1} \implies a_k = \frac{3k-2}{k(3k+5)} a_{k-1}.
          \]
          Resolving this iteratively provides $a_k$ in terms of $a_0$:
          \[
            a_k = a_0 \frac{\prod_{j=1}^k (3j-2)}{k! \prod_{j=1}^k (3j+5)}.
          \]
          The function must satisfy the boundary condition $u(1) = 1$. The series uniformly converges at $x=1$ (ratio test approaches $0$), so we can evaluate it to find $a_0$.
          The unique solution in $X$ is:
          \[
            u(x) = c \cdot x^{1/3} \sum_{k=0}^\infty \frac{\prod_{j=1}^k (3j-2)}{k! \prod_{j=1}^k (3j+5)} x^k \quad \text{where } c = \left( \sum_{k=0}^\infty \frac{\prod_{j=1}^k (3j-2)}{k! \prod_{j=1}^k (3j+5)} \right)^{-1}.
          \]
        \end{enumerate}

  \item % Solution 10
        The function $g(w, u)$ operates by shifting blocks of $\mathbf{a}$'s located after a leading block of $\mathbf{b}$'s into the second argument $u$, simultaneously replacing the first occurrence of $\mathbf{b}\mathbf{a}$ with $\mathbf{a}\mathbf{b}$.
        Let us rigorously trace the evaluation of $f(w) = g(w, \epsilon)$ on a string $w$ that does not consist entirely of $\mathbf{a}$'s followed by $\mathbf{b}$'s. Such a string can be uniquely decomposed as $w = \mathbf{a}^m \mathbf{b}^p \mathbf{a} w'$ where $m \ge 0$ and $p \ge 1$.
        Tracing the recursive rules of $g$:
        \begin{align*}
          f(\mathbf{a}^m \mathbf{b}^p \mathbf{a} w') &= g(\mathbf{a}^m \mathbf{b}^p \mathbf{a} w', \epsilon) \\
          &= g(\mathbf{b}^p \mathbf{a} w', \mathbf{a}^m) \\
          &= g(\mathbf{b}^{p-1} \mathbf{a} w', \mathbf{b} \mathbf{a}^m) \\
          &= g(\mathbf{b} \mathbf{a} w', \mathbf{b}^{p-1} \mathbf{a}^m) \\
          &= \mathbf{b}^{p-1} \mathbf{a}^{m+1} \mathbf{b} w'.
        \end{align*}
        If the string is of the form $w = \mathbf{a}^m \mathbf{b}^p$ (all $\mathbf{a}$'s are packed at the front), the evaluation is:
        \[
          f(\mathbf{a}^m \mathbf{b}^p) = g(\mathbf{b}^p, \mathbf{a}^m) = \dots = g(\mathbf{b}, \mathbf{b}^{p-1} \mathbf{a}^m) = \mathbf{b}^p \mathbf{a}^m.
        \]
        This algorithmic mechanism is exactly the well-known "Revolving Door" or "Lexicographic Next Permutation" generator applied to combinations (multiset permutations of $\{\mathbf{a}^k, \mathbf{b}^{n-k}\}$). 
        The operation $f$ systematically locates the first instance where $\mathbf{b}$ precedes $\mathbf{a}$ ($\mathbf{b}\mathbf{a}$), swaps them to $\mathbf{a}\mathbf{b}$, and consolidates all prior $\mathbf{a}$'s strictly to the right of the newly formed sequence, abutting the swap. 
        This is functionally a bijective successor mapping on the finite set $B_{n,k}$. Starting from the lexicographically "lowest" string $\mathbf{b}^{n-k}\mathbf{a}^k$, repeated applications of $f$ mathematically step through every unique permutation in $B_{n,k}$ exactly once before concluding at the terminal string $\mathbf{a}^k \mathbf{b}^{n-k}$. The mapping $f(\mathbf{a}^k \mathbf{b}^{n-k}) = \mathbf{b}^{n-k} \mathbf{a}^k$ then gracefully wraps the sequence back to the starting string, closing the cycle.
        Since $f$ constitutes a single unbroken permutation cycle of length $|B_{n,k}| = \binom{n}{k}$, its orbit fully generates the set:
        \[
          \{ f^i(\mathbf{b}^{n-k}\mathbf{a}^k) \mid i \in \mathbf{N} \} = B_{n,k}.
        \]
\end{enumerate}

\end{document}
